% !TeX program = pdflatex
% =============================================================================
% Folien: Einstieg zu Stromkreisen (Wiederholung U-I-R)
% UZH Praktikum I, Sara Romer Urban, Kantonsschule im Lee, HS2026
% Klasse: 1f (23.09.2026)
%
% ENTWURF. Einstiegsaufgaben zur Wiederholung des Stoffs aus der
% U-I-R-Lektion (elektrizitaet/u-i-r/gruppenpuzzle3-u-i-r.tex), gedacht als
% Plenumseinstieg vor/zu Beginn der Schaltungslektion, interaktiv an der
% Wandtafel/Leinwand mit der ganzen Klasse:
% (1) Leerer Stromkreis: SuS skizzieren im Heft einen vollstaendigen
%     Stromkreis mit den richtigen Schaltzeichen (Batterie, Widerstand,
%     Schalter, Energiewandler) und sollen Stromrichtung und die Stelle(n)
%     der Spannungsnutzung benennen koennen. Die Reveal-Folie zeigt eine
%     moegliche Loesung.
% (2) Serie-/Parallelschaltung im Bild wiedererkennen: dieselbe Skizze mit
%     zwei Gluehbirnen wie im Gruppenpuzzle (Gruppe C), bewusst OHNE
%     Beschriftung "Serieschaltung"/"Parallelschaltung". SuS ordnen die
%     beiden Schaltungen selbst zu und beschreiben danach rein
%     phaenomenologisch (keine Strom-/Spannungserklaerung, das ist Stoff
%     der naechsten Unterrichtseinheit), wie hell die Gluehbirnen in den
%     beiden Schaltungen im Vergleich brennen.
% (3) Beispielrechnung $U=R\cdot I$ an einem einfachen Stromkreis mit genau
%     einem Widerstand, mit Einheitenpraefixen (kOmega/mA) zur
%     Wiederholung der Praefixumrechnung.
% (4) Richtig-oder-falsch-Runde zu gaengigen Fehlvorstellungen (auf Wunsch):
%     fuenf Aussagen (drei falsche Fehlvorstellungen, zwei richtige zur
%     Gegenprobe, damit "immer falsch" keine gueltige Ratestrategie ist) zu
%     Ladungserzeugung vs. Ladungstrennung in der Batterie,
%     Elektronenwanderung vs. bereits gefuelltem Stromkreis
%     (Wasserkreislaufanalogie), Batterie als Konstantspannungs- statt
%     Konstantstromquelle, und Notwendigkeit eines geschlossenen Kreises.
%     Kein eigenes lernziele.md fuer u-i-r/schaltungen vorhanden, deshalb
%     aus allgemeiner fachdidaktischer Kenntnis gaengiger
%     Schuelervorstellungen zum Stromkreis entwickelt, nicht aus einer
%     Liste bekannter Lernschwierigkeiten uebernommen. Richtig/Falsch wird
%     direkt auf derselben Folie per Beamer-Overlay eingeblendet (ein
%     Haekchen/Kreuz je Klick, kumulativ), keine separaten
%     Auflösungsfolien mit Erklaerungstext.
%
% -----------------------------------------------------------------------------
% Korrektur (23.09.2026, auf Wunsch): Aufgabe 2 "Serie oder Parallel?" war
% hier komplett entfernt worden (siehe vorherige Korrekturnotiz), da der
% urspruengliche Aufgabentext Strom-/Spannungsverhalten erklaeren liess:
% Stoff der naechsten Unterrichtseinheit. Auf Wunsch wieder eingefuegt,
% diesmal aber bewusst rein phaenomenologisch: Die Aufgabe fragt nur noch,
% welche der beiden ungekennzeichneten Skizzen Serie- bzw. Parallelschaltung
% ist (Wiedererkennung im Bild) und wie hell die Gluehbirnen dabei jeweils
% brennen (reine Beobachtung: in der Serieschaltung beide gleich hell, aber
% schwaecher als eine einzelne Gluehbirne allein; in der Parallelschaltung
% beide gleich hell wie eine einzelne Gluehbirne allein), ohne jede
% Erklaerung ueber Ladungstraeger, Strom oder Spannung. Alle nachfolgenden
% Aufgaben entsprechend neu durchnummeriert (Beispielrechnung wieder
% Aufgabe 3, Richtig-oder-falsch-Runde wieder Aufgabe 4).
%
% -----------------------------------------------------------------------------
% Korrektur (23.09.2026, auf Wunsch): Die beiden Auflösungsfolien mit den
% ausformulierten Erklaerungen zur Richtig-oder-falsch-Runde (Aufgabe 4)
% entfernt. Stattdessen zeigt jetzt dieselbe Aussagen-Folie per
% Beamer-Overlay (\only<n->{}) nacheinander ein Haekchen (\ding{51}, gruen)
% oder Kreuz (\ding{55}, rot) neben jeder Aussage, kumulativ mit jedem
% Klick -- dafuer neu \usepackage{pifont} in der Praeambel. Die Folie hat
% dadurch sechs Overlay-Schritte (Ausgangszustand ohne Zeichen, dann ein
% Schritt je Aussage) statt vorher drei separaten Folien.
%
% -----------------------------------------------------------------------------
% Korrektur (23.09.2026, auf Wunsch): (1) Der Begriff "Verbraucher" im
% ganzen Dokument fallengelassen, durchgaengig "Wandler" verwendet (Aufgabe
% 1: Bauteilliste "Energiewandler" ohne die Verbraucherklammer mehr,
% Schaltplanbeschriftung "Wandler" statt "Verbraucher", Aussage 5 in
% Aufgabe 4 ebenso), konsistent mit dem in arbeitsblatt4-schaltungen.tex
% bereits etablierten Begriff "Wandler" (Merksatz Serieschaltung: "die
% Bauteile, die Wandler"). (2) Neue sechste Aussage in der
% Richtig-oder-falsch-Runde (Aufgabe 4, auf Wunsch): "Eine Gluehbirne
% verbraucht Energie." Falsch, analog zur bereits an anderer Stelle
% behandelten "Stromverbrauch"-Fehlvorstellung: Die Gluehbirne wandelt
% elektrische Energie um (in Licht und Waerme), sie verbraucht sie nicht
% im Sinne einer Vernichtung: Energie bleibt erhalten. Overlay-Schritte
% der Folie entsprechend von sechs auf sieben erweitert.
% =============================================================================
\documentclass[14pt,aspectratio=169]{beamer}

\usepackage[T1]{fontenc}
\usepackage[utf8]{inputenc}
\usepackage[ngerman]{babel}
\usepackage{amsmath}
\usepackage{siunitx}
% Hausstil: zusammengesetzte Einheiten immer als echter Bruch (\frac), nicht
% als Inline-Schraegstrich oder \tfrac -- gilt global fuer jedes \si/\SI.
\sisetup{per-mode=fraction,fraction-command=\frac}
\usepackage{tikz}
\usepackage{physics}
\usepackage[european]{circuitikz}
\usepackage{enumitem}
\setlist[itemize]{itemsep=0.2em,topsep=0.3em,label=\textbullet}
\usepackage{pifont}

% --- KFR-Theme: minimal, hoher Kontrast, grosse Schrift ---------------------
\usetheme{default}
\usefonttheme{professionalfonts}
\setbeamertemplate{navigation symbols}{}
\setbeamertemplate{sidebar}{}
\setbeamertemplate{headline}{}
\beamertemplatenavigationsymbolsempty

\definecolor{kfrblau}{RGB}{20,58,94}
\definecolor{kfrgrau}{RGB}{90,90,90}

\setbeamercolor{normal text}{fg=black,bg=white}
\setbeamercolor{structure}{fg=kfrblau}
\setbeamercolor{title}{fg=kfrblau,bg=white}
\setbeamercolor{frametitle}{fg=white,bg=kfrblau}
\setbeamercolor{footline}{fg=kfrgrau,bg=white}
\setbeamercolor{page number in head/foot}{fg=kfrgrau}
\setbeamercolor{itemize item}{fg=kfrblau}
\setbeamercolor{itemize subitem}{fg=kfrblau}

\setbeamerfont{frametitle}{size=\Large,series=\bfseries}
\setbeamerfont{title}{size=\LARGE,series=\bfseries}
\setbeamerfont{itemize item}{size=\normalsize}

\setbeamertemplate{itemize item}{\textbullet}
\setbeamertemplate{itemize subitem}{--}

% Schlichte Fusszeile: nur Seitenzahl, kein Autoren-/Konferenz-Ballast
\setbeamertemplate{footline}{%
  \hfill%
  \usebeamercolor[fg]{page number in head/foot}%
  \usebeamerfont{page number in head/foot}%
  \insertframenumber\,/\,\inserttotalframenumber\hspace{1em}\vspace{0.5em}
}

% Schlichter Frametitle-Balken (voller Breite, kein Untertitel-Chrome)
\setbeamertemplate{frametitle}{%
  \nointerlineskip
  \begin{beamercolorbox}[wd=\paperwidth,ht=1.4em,dp=0.4em,leftskip=0.5em]{frametitle}
    \usebeamerfont{frametitle}\insertframetitle
  \end{beamercolorbox}
}

\title{Stromkreise}
\subtitle{Einstieg: Wiederholung U-I-R}
\author{}
\institute{UZH Praktikum I, Klasse 1f}
\date{23.09.2026}

\begin{document}

% --- Titelfolie --------------------------------------------------------------
\begin{frame}[plain]
  \centering
  {\usebeamerfont{title}\usebeamercolor[fg]{title}\inserttitle\par}
  \vspace{0.3em}
  {\large\insertsubtitle\par}
  \vspace{1.5em}
  {\normalsize\insertinstitute, \insertdate\par}
\end{frame}

% =============================================================================
% Aufgabe 1: Leerer Stromkreis
% =============================================================================
\begin{frame}{Aufgabe 1: Stromkreis skizzieren}
  \centering
  \begin{circuitikz}[scale=0.8]
    \draw (0,0) -- (0,2.2) -- (5,2.2) -- (5,0) -- (0,0);
  \end{circuitikz}

  \begin{minipage}{0.92\textwidth}
  \footnotesize
  Skizzieren Sie in Ihrem Heft einen vollständigen Stromkreis. Verwenden
  Sie dabei die richtigen Schaltzeichen für:
  \begin{itemize}
    \item Batterie (Spannungsquelle)
    \item Widerstand
    \item Schalter
    \item Energiewandler
  \end{itemize}
  Zeichnen Sie ausserdem ein: \textbf{Wo und in welche Richtung fliesst
  der Strom?} \textbf{Wo wird eine Spannung benötigt?}
  \end{minipage}
\end{frame}

\begin{frame}{Aufgabe 1: Mögliche Lösung}
  \centering
  \begin{circuitikz}[scale=0.75]
    \draw (0,0) to[battery1,invert,l=$U$] (0,2.2)
                to[normal open switch] (1.8,2.2)
                to[R,l=$R$] (4.1,2.2)
                to[lamp,l=Wandler] (6.4,2.2)
                -- (6.4,0)
                to[short,i>=$I$] (0,0);
  \end{circuitikz}

  \begin{minipage}{0.92\textwidth}
  \footnotesize
  Es gibt viele richtige Lösungen. Wichtig ist nur:
  \begin{itemize}
    \item Der Strom $I$ fliesst (in der technischen Stromrichtung) vom
    Pluspol der Batterie durch den ganzen Stromkreis zum Minuspol zurück.
    \item Eine Spannung wird dort benötigt, wo elektrische Energie
    umgewandelt wird: am Widerstand und am Energiewandler. Am Schalter
    und an den Kabeln selbst fällt (bei geschlossenem Schalter) keine
    Spannung ab.
  \end{itemize}
  \end{minipage}
\end{frame}

% =============================================================================
% Aufgabe 2: Serie oder Parallel? (rein phaenomenologisch)
% =============================================================================
\begin{frame}{Aufgabe 2: Serie oder Parallel?}
  \centering
  \begin{tabular}{cc}
  \begin{circuitikz}[scale=0.8]
    \draw (0,0) to[battery1,invert,l=$U$] (0,3.4) -- (4,3.4)
                to[lamp] (4,2.1)
                -- (4,1.3)
                to[lamp] (4,0)
                -- (0,0);
  \end{circuitikz}
  &
  \begin{circuitikz}[scale=0.8]
    \draw (0,0) to[battery1,invert,l=$U$] (0,3);
    \draw (0,3) -- (4,3);
    \draw (0,0) -- (4,0);
    \draw (2,3) to[lamp] (2,0);
    \draw (4,3) to[lamp] (4,0);
  \end{circuitikz}
  \\
  {\small Schaltung A} & {\small Schaltung B}
  \end{tabular}

  \vspace{0.8em}
  \begin{minipage}{0.85\textwidth}
  \small
  Welche der beiden Schaltungen ist eine \textbf{Serieschaltung}, welche
  eine \textbf{Parallelschaltung}? Woran erkennen Sie das im Bild?

  Stellen Sie sich vor, Sie bauen beide Schaltungen tatsächlich auf: Wie
  hell brennen die Glühbirnen jeweils?
  \end{minipage}
\end{frame}

\begin{frame}{Aufgabe 2: Auflösung}
  \centering
  \begin{tabular}{cc}
  \begin{circuitikz}[scale=0.75]
    \draw (0,0) to[battery1,invert,l=$U$] (0,3.4) -- (4,3.4)
                to[lamp] (4,2.1)
                -- (4,1.3)
                to[lamp] (4,0)
                -- (0,0);
  \end{circuitikz}
  &
  \begin{circuitikz}[scale=0.75]
    \draw (0,0) to[battery1,invert,l=$U$] (0,3);
    \draw (0,3) -- (4,3);
    \draw (0,0) -- (4,0);
    \draw (2,3) to[lamp] (2,0);
    \draw (4,3) to[lamp] (4,0);
  \end{circuitikz}
  \\
  {\small A: Serieschaltung} & {\small B: Parallelschaltung}
  \end{tabular}

  \vspace{0.6em}
  \begin{minipage}{0.92\textwidth}
  \small
  \textbf{Serieschaltung (A):} Beide Glühbirnen brennen gleich hell wie
  einander, aber schwächer als eine einzelne Glühbirne allein an
  derselben Batterie.

  \textbf{Parallelschaltung (B):} Beide Glühbirnen brennen gleich hell
  wie einander, und zwar genauso hell wie eine einzelne Glühbirne allein
  an derselben Batterie.
  \end{minipage}
\end{frame}

% =============================================================================
% Aufgabe 3: Beispielrechnung U = R * I
% =============================================================================
\begin{frame}{Aufgabe 3: Beispielrechnung}
  \centering
  \begin{circuitikz}[scale=1.15]
    \draw (0,0) to[battery1,invert,l=$U$] (0,3)
                to[R,l=$R$] (5,3)
                -- (5,0)
                to[short,i>=$I$] (0,0);
  \end{circuitikz}

  \vspace{0.8em}
  \begin{minipage}{0.8\textwidth}
  \small
  Gegeben: $R=\SI{1.5}{\kilo\ohm}$, $I=\SI{4}{\milli\ampere}$.

  Gesucht: die Spannung $U$.
  \end{minipage}
\end{frame}

\begin{frame}{Aufgabe 3: Lösung}
  \small
  Formel:
  \[
    U = R\cdot I
  \]
  Einheiten umrechnen (Kilo $=10^3$, Milli $=10^{-3}$):
  \[
    R = \SI{1.5}{\kilo\ohm} = 1.5\times10^{3}\,\Omega = \SI{1500}{\ohm}
  \]
  \[
    I = \SI{4}{\milli\ampere} = 4\times10^{-3}\,\text{A} = \SI{0.004}{\ampere}
  \]
  Einsetzen:
  \[
    U = \SI{1500}{\ohm}\cdot\SI{0.004}{\ampere}
  \]
  Ausrechnen:
  \[
    U = \SI{6}{\volt}
  \]
\end{frame}

% =============================================================================
% Aufgabe 4: Richtig oder falsch? (Fehlvorstellungen)
% =============================================================================
\begin{frame}{Aufgabe 4: Richtig oder falsch?}
  \footnotesize
  \begin{itemize}[leftmargin=1.6em]
    \item[1.] Die Batterie erzeugt die Elektronen, die als Strom fliessen.
    \only<2->{\textcolor{red!70!black}{\textbf{\ding{55}}}}
    \item[2.] Wird der Schalter geschlossen, müssen die Elektronen zuerst
    von der Batterie zur Glühbirne wandern, bevor diese zu leuchten
    beginnt.
    \only<3->{\textcolor{red!70!black}{\textbf{\ding{55}}}}
    \item[3.] Ein Schalter unterbricht den Stromkreis unabhängig davon,
    an welcher Stelle der Schleife er eingebaut ist.
    \only<4->{\textcolor{green!50!black}{\textbf{\ding{51}}}}
    \item[4.] Die Batterie liefert eine konstante Stromstärke.
    \only<5->{\textcolor{red!70!black}{\textbf{\ding{55}}}}
    \item[5.] Damit Strom fliesst, genügt ein einzelnes Kabel von der
    Batterie zum Wandler.
    \only<6->{\textcolor{red!70!black}{\textbf{\ding{55}}}}
    \item[6.] Eine Glühbirne verbraucht Energie.
    \only<7->{\textcolor{red!70!black}{\textbf{\ding{55}}}}
  \end{itemize}
\end{frame}

\end{document}
